%%  THE GENERATIVE-AI DECLARATION.  Elsevier's policy (updated June 2026,
%%  read 2026-09-24) asks for it "at the end of your manuscript,
%%  immediately before the references", under the section title below, and
%%  asks that AI-written code be described in the Methods section, which
%%  Section 4.1 does.  It names no author, so the blind article carries it
%%  (90_backmatter.tex) and the title page repeats it (assemble_paper.py).
%%  The title stays on ONE source line: texlint records each section's name
%%  into results/section_words.csv, and a wrapped title broke that CSV.
\section*{Declaration of generative AI and AI-assisted technologies in the manuscript preparation process}

During the preparation of this work the \ifblind author(s)\else author\fi\
used Anthropic Claude (\aiModelId; \aiModelDate) in order to draft and edit
prose, to write analysis and verification code, and to conduct internal
critical reads of the manuscript against its code and results. After using this tool
the \ifblind author(s)\else author\fi\ reviewed and edited the content as
needed and \ifblind take(s)\else takes\fi\ full responsibility for the
content of the publication.

The use in writing code is described in the Methods section
(\ifdefined\titlepagedocument the article's Section~4.1\else
Section~\ref{sec:boot}\fi). No
generative model produced, imputed, augmented or selected any datum, any
result or any citation behind a claim of this paper: every result is
produced by archived code that one command re-runs, and every reference was
checked against its published record. Every session's date, and its model
identifier where one was recorded, are in \texttt{AI-USE.md} in the code
repository. The one
place where a machine-assisted judgement enters a study is the literature
pilot of Supplement~\ref{app:pilot}, which is declared where it is
described, and no claim in this paper rests on it.
